\begin{lemma}
Let $g\in\mathsf{IncSeq}$ be nonidentity.  There is a prefix-continuous map
$\rho:\mathsf{IncSeq}\to\mathsf{IncSeq}$ satisfying
\[
\rho(f\circ g)=\rho(f)\circ\mathsf{SuccSeq}
\]
for every $f$.  Equivalently, right composition by $g$ continuously
homomorphically maps to the ordinary shift.
\end{lemma}
\begin{proof}
By the nonidentity hypothesis and \noderef{FirstMovedPoint}, choose the
least moved point $k_g$ with $k_g<g(k_g)$.  Apply \noderef{OrbitEmbedding} to
this inequality and
choose $G\in\mathsf{IncSeq}$ satisfying
\[
G(n)=\mathsf{OrbitPoint}(g,k_g,n)
\]
for every $n$.  By the definition of \noderef{OrbitPoint}, this means
$G(n)=g^{[n]}(k_g)$.  For $f\in\mathsf{IncSeq}$ define
\[
\rho(f)=f\circ G.
\]
This value is in $\mathsf{IncSeq}$: by \noderef{IncSeq}, both $f$ and $G$
are order embeddings, and their composite is order preserving and
injective, hence is again an order embedding.

We verify prefix-continuity using the criterion in \noderef{BaireContinuous}
and the meaning of prefix agreement in \noderef{PrefixAgree}.  Fix
$f\in\mathsf{IncSeq}$ and an output length $m$.  If $m=0$, choose input
length $k=0$; the conclusion is the vacuous statement
\(\operatorname{PrefixAgree}(0,\rho(f),\rho(h))\) for every $h$.
Suppose instead that $m>0$, and choose
\[
k=G(m-1)+1.
\]
Let $h\in\mathsf{IncSeq}$ satisfy
$\operatorname{PrefixAgree}(k,f,h)$.  For every $i<m$, we have
$i\leq m-1$, and the order-preserving property of the order embedding
$G$ gives $G(i)\leq G(m-1)<k$.  Thus prefix agreement gives
$h(G(i))=f(G(i))$.  Consequently
\[
\rho(h)(i)=h(G(i))=f(G(i))=\rho(f)(i)
\]
for every $i<m$, so
$\operatorname{PrefixAgree}(m,\rho(f),\rho(h))$.  This supplies the
witness required by \noderef{BaireContinuous} for every $f$ and $m$.

It remains to prove the homomorphism equation.  From
\noderef{OrbitPoint} and the displayed choice of $G$, the successor rule
for function iterates gives, for every $n$,
\[
g(G(n))=g\bigl(g^{[n]}(k_g)\bigr)=g^{[n+1]}(k_g)=G(n+1).
\]
This calculation uses only the definition of the orbit and the iterate
recursion, not the proof of \noderef{OrbitEmbedding}.  Therefore, for every
$f\in\mathsf{IncSeq}$ and $n\in\mathbb N$,
\[
\begin{aligned}
\rho(f\circ g)(n)
  &=f\bigl(g(G(n))\bigr)\\
  &=f(G(n+1))\\
  &=\rho(f)(n+1).
\end{aligned}
\]
By \noderef{RightComp}, $f\circ g=\mathsf{RightComp}(g)(f)$.  By
\noderef{ShiftMap} and \noderef{SuccSeq},
\[
\rho(f)\circ\mathsf{SuccSeq}=\mathsf{ShiftMap}(\rho(f)),
\qquad
(\rho(f)\circ\mathsf{SuccSeq})(n)=\rho(f)(n+1).
\]
Thus the coordinate equalities give
\[
\rho(\mathsf{RightComp}(g)(f))
 =\mathsf{ShiftMap}(\rho(f))
\]
for every $f$.  Together with the prefix-continuity just proved, this is a
\noderef{ContMor} witness, and hence establishes the stated conclusion.
\end{proof}
